\section{研究背景与实验目标}
卫星导航在室内环境中会受到遮挡、反射与衰减影响，难以持续提供稳定的米级定位。课程内容从定位传感器、几何定位、航迹推算、数据库匹配扩展到运动约束、多源融合和机器学习。本实验选择项目中已有数据和程序覆盖的三条技术路线，重点考察它们的观测量、误差传播与适用条件，而不是把所有课件主题都误当作已完成的实验。

实验目标是：（1）从代码和数据字段还原真实算法，而非仅按目录名称推断；（2）对具有独立输入的数据复现计算，或对现有程序输出进行可追溯复核；（3）以参考位置评价定位误差，并区分测距残差、轨迹末点偏差等不同指标；（4）指出实验的可重复性、坐标对准和参数标定限制。

\section{环境、数据与证据范围}
项目代码使用 C++17 实现距离交会和 PDR，Python 实现 RSS 指纹及传感器绘图/转换。计算与绘图脚本运行于 Python 3 环境，依赖 NumPy、Matplotlib；C++ 工程 README 指定支持 C++17 的编译器。当前报告的独立复算脚本位于 \texttt{analysis.py}，所有本报告新增图表和统计文件保存在本目录，原始数据与代码未修改。

\begin{table}[htbp]
\centering
\caption{数据与本报告分析用途}
\label{tab:data}
\begin{tabular}{p{3.0cm}p{5.3cm}p{5.2cm}}
\toprule
数据集 & 结构与观测 & 本报告中的用途及限制 \\
\midrule
距离交会：教室超声 & 基站坐标、逐历元测距、逐历元参考位置（CSV/文本） & 复核项目已有 C++ 位置输出；参考位置按时间线性插值，评价二维误差。输出文件是项目既有结果，本次未在原目录重新运行 C++。 \\
RSS 指纹 & 97 个参考点坐标、1,989 条建库 RSS 记录、50 条带参考坐标测试记录；每条包含 30 个 AP 通道，缺失值为 $-110$ dBm & 按现有 Python 算法在独立脚本中重新建库、定位和统计误差。 \\
PDR & 6,829 行、约 0.05 s 间隔的 12 列传感器矩阵；588 个带时间戳的参考经纬度位置 & 对项目 C++ 已保存的定位轨迹进行独立对齐和误差统计；不声称本轮重新编译运行 PDR C++。 \\
Android 采集样例 & 采集 App 的传感器、Wi-Fi 和 GPS/KML 记录，包含陀螺、加速度、磁场、Wi-Fi 扫描等 & 用于说明采集链和数据转换用途。该样例不等同于本报告指纹库/PDR 数据，未把它混入定位精度统计。 \\
无人机、UWB、学生超声数据 & 多组基站和测距文件，部分带高度、真值或遮挡分组 & 代码包含对应处理程序；本报告不报告这些组的精度，因为本次没有在隔离副本中运行其 C++ 主程序并完成真值口径核验。 \\
\bottomrule
\end{tabular}
\end{table}

指纹训练文件按参考点编号和 30 个 RSS 数值组织，测试文件另含时间与三维参考坐标。PDR 的传感器矩阵列顺序由读取实现确认：时间、三轴角速度、三轴加速度、三轴磁场以及两个未参与定位的字段。代码将角速度从度每秒转为弧度每秒，将磁场数值乘 $10^{-7}$ 转为特斯拉。距离交会教室程序读取四个基站以及各历元测距/参考数据；输出字段中的 \texttt{range\_rms\_m} 表示测距拟合残差。
